\begin{abstract}
Surprising videos, \emph{e.g.}, funny clips, creative performances, or visual illusions, attract significant attention. Enjoyment of these videos is not simply a response to visual stimuli; rather, it hinges on the human capacity to understand (and appreciate) commonsense \emph{violations} depicted in these videos.
We introduce \textbf{FunQA}, a challenging video question answering (QA) dataset specifically designed to evaluate and enhance the depth of video reasoning based on counter-intuitive and fun videos. Unlike most video QA benchmarks which focus on less surprising contexts, e.g., cooking or instructional videos, FunQA covers three previously unexplored types of surprising videos: \emph{1) HumorQA}, \emph{2) CreativeQA}, and \emph{3) MagicQA}.
For each subset, we establish rigorous QA tasks designed to assess the model's capability in counter-intuitive timestamp localization, detailed video description, and reasoning around counter-intuitiveness.
We also pose higher-level tasks, such as attributing a fitting and vivid title to the video, and scoring the video creativity.
In total, the FunQA benchmark consists of 312K free-text QA pairs derived from 4.3K video clips, spanning a total of 24 video hours.
Extensive experiments with existing VideoQA models reveal significant performance gaps for the FunQA videos across spatial-temporal reasoning, visual-centered reasoning, and free-text generation.
 \end{abstract}

